\chapter{The Canonical Equations}

\textbooknumber{\S 40}
\section{Hamilton's equations}

% ============ Problem 1 ============ %

\begin{problem}
{
Find the Hamiltonian for a single particle in Cartesian, cylindrical and spherical co-ordinates.
}
{}
{}
\end{problem}

% ============ Problem 2 ============ %

\begin{problem}
{
Find the Hamiltonian for a particle in a uniformly rotating frame of reference.
}
{}
{}
\end{problem}

% ============ Problem 3 ============ %

\begin{problem}
{
Find the Hamiltonian for a system comprising one particle of mass $M$ and $n$ particles each of mass $m$, excluding the motion of the centre of mass (see \probref{\S 13}{1}).
}
{}
{}
\end{problem}

\textbooknumber{\S 41}
\section{The Routhian}

% ============ Problem 1 ============ %

\begin{problem}
{
Find the Routhian for a symmetrical top in an external field $U(\phi, \theta)$, eliminating the cyclic co-ordinate $\psi$ (where $\psi$, $\phi$, $\theta$ are Eulerian angles).
}
{}
{}
\end{problem}

\textbooknumber{\S 42}
\section{Poisson brackets}

% ============ Problem 1 ============ %

\begin{problem}
{
Determine the Poisson brackets formed from the Cartesian components of the momentum $\mathbf{p}$ and the angular momentum $\mathbf{M} = \mathbf{r} \times \mathbf{p}$ of a particle.
}
{}
{}
\end{problem}

% ============ Problem 2 ============ %

\begin{problem}
{
Determine the Poisson brackets formed from the components of $\mathbf{M}$.
}
{}
{}
\end{problem}

% ============ Problem 3 ============ %

\begin{problem}
{
Show that $[\phi, M_z] = 0$, where $\phi$ is any function, spherically symmetrical about the origin, of the co-ordinates and momentum of a particle.
}
{}
{}
\end{problem}

% ============ Problem 4 ============ %

\begin{problem}
{
Show that $[\mathbf{f}, M_z] = \mathbf{f} \times \mathbf{n}$, where $\mathbf{f}$ is a vector function of the co-ordinates and momentum of a particle, and $\mathbf{n}$ is a unit vector parallel to the $z$-axis.
}
{}
{}
\end{problem}

\textbooknumber{\S 44}
\section{Maupertuis' principle}

% ============ Problem 1 ============ %

\begin{problem}
{
Derive the differential equation of the path from the variational principle (44.10).
}
{}
{}
\end{problem}

\textbooknumber{\S 48}
\section{Separation of the variables}

% ============ Problem 1 ============ %

\begin{problem}
{
Find a complete integral of the Hamilton--Jacobi equation for motion of a particle in a field $U = \alpha/r - Fz$ (a combination of a Coulomb field and a uniform field), and find a conserved function of the co-ordinates and momenta that is specific to this motion.
}
{}
{}
\end{problem}

% ============ Problem 2 ============ %

\begin{problem}
{
The same as \probref{\S 48}{1}, but for a field $U = \alpha_1/r_1 + \alpha_2/r_2$ (the Coulomb field of two fixed points at a distance $2\sigma$ apart).
}
{}
{}
\end{problem}

\textbooknumber{\S 50}
\section{Canonical variables}

% ============ Problem 1 ============ %

\begin{problem}
{
Write down the equations of motion in canonical variables for a harmonic oscillator (whose Hamiltonian is (49.11)) with time-dependent frequency.
}
{}
{}
\end{problem}

\textbooknumber{\S 51}
\section{Accuracy of conservation of the adiabatic invariant}

% ============ Problem 1 ============ %

\begin{problem}
{
Estimate $\delta I$ for a harmonic oscillator with a frequency that varies slowly according to
\begin{equation*}
    \omega^2 = \omega_0^2 \, \frac{1 + a e^{\alpha t}}{1 + e^{\alpha t}}
\end{equation*}
from $\omega_- = \omega_0$ for $t = -\infty$ to $\omega_+ = \sqrt{a}\,\omega_0$ for $t = +\infty$ ($a > 0$, $\alpha \ll \omega_0$).
}
{}
{}
\end{problem}

% ============ Problem 2 ============ %

\begin{problem}
{
A particle oscillates in a potential well. Determine how its energy varies under a frictional force $f_{\text{fr}} = -\alpha \Dot{x}$ with a small coefficient $\alpha$ ($x$ being a Cartesian co-ordinate).
}
{}
{}
\end{problem}

\textbooknumber{\S 52}
\section{Conditionally periodic motion}

% ============ Problem 1 ============ %

\begin{problem}
{
Calculate the action variables for elliptic motion in a field $U = -\alpha/r$.
}
{}
{}
\end{problem}
